\documentclass[10pt,a4paper]{amsart}
\usepackage[margin=19mm]{geometry}
\usepackage{amsmath,amssymb,mathtools}
\usepackage[scr=rsfs,cal=euler]{mathalfa}
\usepackage{xfrac}
\usepackage[unicode,hidelinks,bookmarks=false]{hyperref}
\newcommand{\ie}{i.e.\ }
\newcommand{\cf}{cf.\ }
\newcommand{\resp}{respectively\ }
\newcommand{\Subsection}{\subsection}
\providecommand{\nobreakdash}{\nobreak\hbox{-}}
\setlength{\emergencystretch}{2em}
\multlinegap0pt
\usepackage{mathtools}
\usepackage{amssymb}
\usepackage[shortlabels]{enumitem}
\usepackage[normalem]{ulem}
\newtheorem{lemma}{Lemma}
\newcommand{\K}{\mathbb{K}}
\newcommand{\conv}{\ast}
\newcommand{\kprod}{\otimes}
\title{The Geometry of Polynomial Group Convolutional Neural Networks}
\author{Yacoub Hendi, Daniel Persson, Magdalena Larfors}
\date{Source: arXiv:2603.29566v3 (CC BY 4.0)}
\begin{document}
\maketitle

\noindent\emph{Source and licence.} The Geometry of Polynomial Group Convolutional Neural Networks. Version arXiv:2603.29566v3 (CC BY 4.0), distributed by arXiv under the Creative Commons Attribution 4.0 International licence, http://creativecommons.org/licenses/by/4.0/, as recorded in the arXiv licence metadata for this exact version and shown on the abstract page https://arxiv.org/abs/2603.29566v3. Full licence notice: \texttt{licenses/arxiv-2603.29566-CC-BY-4.0.txt}.\par\smallskip

\noindent\textit{Editorial scope.} Counted unit: the article's lemma jac\_varphi\_simplified\_case (source0.tex lines 1107--1109). The article's complete original proof of it is delivered with it (source0.tex lines 1110--1173, one \texttt{proof} environment of 64 lines). No mathematics is written, repaired or re-derived: the statement and the proof are the article's own lines, in the article's own notation, and no retained line is substituted or re-flowed.  Retained uncounted as the source-local context the proof needs: lines 1066--1086.  Nothing was omitted from any retained span, so the package carries no omission record.  No retained line contains a citation, so the wrapper carries no bibliography and no bibliography entry.  No retained line contains a figure environment, an image-inclusion command or drawing code, so no image is delivered and the package declares no asset dependency.  Editorial formatting only, in addition: an AMS \texttt{amsart} preamble that restates the article's own declarations and macros used by the retained lines, namely the environments \texttt{lemma} on separate counters, together with the standard packages those lines need.  The retained environments print in this document as Lemma 1 in order of appearance, whereas the article prints the same items with the numbering of its own full document.  The complete native source of the article as distributed in the arXiv e-print for this exact version is bundled unchanged as \texttt{sources/197524/source0.tex}.
\begin{lemma}\label{lemma: jacobian_rules_conv_kprod}
    Let $\theta$ and $\psi$ be $G$-filters, then
    \begin{enumerate}
        \item   
        \[
    J_{\theta , \psi}(\theta \conv \psi)(\dot{\theta}, \dot{\psi}) = \dot{\theta}\conv \psi + \theta\conv \dot{\psi} = [\dot{\theta}\conv \psi].
    \]
    \item
\begin{align*}
        J_{\theta}(\theta^{\otimes r})(\dot{\theta}) &= \dot{\theta}\kprod \theta \kprod\hdots\kprod \theta + \theta \kprod \dot{\theta}  \kprod\hdots\kprod \theta + \hdots + \theta \kprod \hdots \kprod \theta \kprod \dot{\theta}\\
        &= \dot{\theta} \kprod \theta^{\otimes (r-1)} + \theta \kprod \dot{\theta} \kprod \theta^{\otimes (r-2)} + \hdots + \theta^{\otimes (r-1)} \kprod \dot{\theta}\\
        &= [\dot{\theta}\kprod \theta \kprod\hdots\kprod \theta].
\end{align*}

    \item Let $\theta = (\theta_1, \hdots, \theta_L)$, $\dot{\theta} = (\dot{\theta}_1, \hdots, \dot{\theta}_L)$. Moreover, let $\theta'$ and $\dot{\theta}'$ be the tuples containing the first $L-1$ filters in $\theta$ and $\dot{\theta}$, respectively. 
    Since $\varphi_\theta = \varphi_{\theta'}^{\otimes r} \conv \theta_L|^{G^{r^{L-1}}}$, we get
    \[
    J_\theta\varphi(\dot{\theta}) = [J_{\theta'}\varphi(\dot{\theta}') \kprod \underbrace{\varphi_{\theta'} \kprod \hdots \kprod \varphi_{\theta'}}_{= \varphi_{\theta'}^{\otimes (r-1)}}] \conv \theta_L + \varphi_{\theta'}^{\otimes r} \conv \dot{\theta}_L.
    \]
    \end{enumerate}
\end{lemma}

\begin{lemma}\label{lemma: jac_varphi_simplified_case}
    Let $\theta = (\theta_1, \hdots, \theta_{L-1}, e) \in \K[G]^L$ where $\theta_i$ is a general filter for $1\leq i < L$ and $e \in G \subset \K[G]$ is the identity element. If $\dot \theta = (\dot \theta_1, \dots, \dot \theta_L) \in \ker J_\theta \varphi$ then $\dot \theta_L = \lambda e$ for some $\lambda \in \K$.
\end{lemma}

\begin{proof}
Write $\theta'=(\theta_1,\dots,\theta_{L-1})$ and similarly $\dot\theta'=(\dot\theta_1,\dots,\dot\theta_{L-1})$.
For brevity, write,
\[
J := J_{\theta'}\varphi(\dot{\theta}').
\]
By Lemma \ref{lemma: jacobian_rules_conv_kprod} (the Jacobian computation rules),
$\dot\theta\in\ker\bigl(J_{\theta}\varphi\bigr)$ if and only if

\begin{equation}\label{eq:jac_kernel_at_theta_L_equal_e}
\bigl[J\kprod \varphi_{\theta'}^{\otimes (r-1)}\bigr]
\;=\;
-\,\varphi_{\theta'}^{\otimes r}\conv \dot\theta_L
\end{equation}
Evaluate \eqref{eq:jac_kernel_at_theta_L_equal_e} at tuples of the special form $\tilde g_i=(g_i,\dots,g_i)\in G^{r^{L-1}}$
(where $g_i\in G^{r^{L-2}}$). This yields, for each $i$,
\[
r\,(\varphi_{\theta'}(g_i))^{r-1} J(g_i) \;=\; \sum_{h\in G} (\varphi_{\theta'}(g_i h^{-1}))^r \dot\theta_L(h).
\]
Hence
\begin{equation}\label{eq:J_at_gi}
J(g_i)
\;=\;
\frac{1}{r}\,\varphi_{\theta'}(g_i)\,\dot\theta_L(e)
\;+\;
\frac{1}{r}\sum_{h\in G\setminus\{e\}}
\frac{(\varphi_{\theta'}(g_i   h^{-1}))^r}{(\varphi_{\theta'}(g_i))^{\,{r-1}}}\,\dot\theta_L(h).
\end{equation}

Now evaluate \eqref{eq:jac_kernel_at_theta_L_equal_e} at $\tilde g=(g_{i_1},\dots,g_{i_r})$ to get,
\[
\sum_{k=1}^r J(g_{i_k}) \prod_{j\neq k} \varphi_{\theta'}(g_{i_j}) = \sum_{h\in G}\varphi_{\theta'}^{\otimes r}(\tilde{g}   h^{-1})\dot{\theta}_L(h).
\]
Substituting \eqref{eq:J_at_gi} in the above equation gives a linear system in the unknowns $\{\dot\theta_L(h)\}_{h\neq e}$ with coefficient matrix $C$ whose entries are symbolic expressions in the values $\varphi_{\theta'}(g_i)$. Denote these coefficients by 
\[
c_{\tilde{g}, h} = \dfrac{1}{r} \sum_{k=1}^r \left(\dfrac{(\varphi_{\theta'}(g_{i_k}   h^{-1}))^r\prod_{j\neq k}\varphi_{\theta'}(g_{i_j})}{(\varphi_{\theta'}(g_{i_k}))^{r-1}}\right) - \varphi_{\theta'}^{\otimes r}(\tilde{g} h^{-1}).
\]

Observe $c_{\tilde g,e}=0$ for every $\tilde g$ that is why we exclude the variable $\dot{\theta}_L(e)$ from the unknowns. The matrix $C$ has $|G^{r^{L-1}}|=n^{r^{L-1}}$ rows and $n-1$ columns, so $C$ is full column rank if and only if some $(n-1)\times(n-1)$ minor is nonzero as a symbolic polynomial in the variables $\theta'$.

Choose the $(n-1)$ rows corresponding to
\[
\tilde g^{\,i}:=(g_1, g_i,\dots,g_i)\qquad(2\le i\le n),
\]
where $g_j=(h_j,\dots,h_j)\in G^{r^{L-2}}$, $h_1=e$ and $h_j \in G$. For this choice, one checks that the diagonal entries of the resulting $(n-1)$-minor contain a factor $(\varphi_{\theta'}(g_1))^{r+1}$ which does not appear elsewhere in that minor. 
Namely,
 \[
 c_{\tilde{g}^i, h_j} = \dfrac{1}{r}\dfrac{(\varphi_{\theta'}(g_i))^{r-1}}{(\varphi_{\theta'}(g_1))^{r-1}} \varphi_{\theta'}((g_1    h_j^{-1}))^r 
 + 
 \dfrac{r-1}{r}\dfrac{(\varphi_{\theta'}(g_i   h^{-j}))^r}{\varphi_{\theta'}(g_i)} \varphi_{\theta'}(g_1) 
 - 
 \varphi_{\theta'}^{\otimes r}(\tilde{g}^{i}    h_j^{-1}).
 \]
and 
 \[
 c_{\tilde{g}^i, h_i} = \dfrac{1}{r}\dfrac{(\varphi_{\theta'}(g_i))^{r-1}}{(\varphi_{\theta'}(g_1))^{r-1}} \varphi_{\theta'}((g_1    h_i^{-1}))^r 
 + 
 \dfrac{r-1}{r}\dfrac{(\varphi_{\theta'}(g_1))^{r+1}}{\varphi_{\theta'}(g_i)}
 - 
 \varphi_{\theta'}^{\otimes r}(\tilde{g}^{i}    h_i^{-1}).
 \]

Since $\varphi_{\theta'}(g_1)$ has monomials in the parameters of $\theta'$ that do not occur  in  $\varphi_{\theta'}(g_i)$ for $2 \geq i \leq n$, then this symbolic minor is nonvanishing for general $\theta'$, so $C$ is symbolically full rank. Therefore the only solution is $\dot\theta_L(h)=0$ for all $h\ne e$. 
\end{proof}

\end{document}
